%==============================================================
%  backmatter/glossaire.tex — Glossaire mathématique
%==============================================================

\newglossaryentry{suite}{
  name={suite},
  description={Application de $\mathbb{N}$ dans un ensemble $E$, notée $(u_n)_{n\in\mathbb{N}}$}
}

\newglossaryentry{suite-monotone}{
  name={suite monotone},
  description={Suite croissante ($u_{n+1}\geq u_n$ pour tout $n$) ou décroissante ($u_{n+1}\leq u_n$ pour tout $n$)}
}

\newglossaryentry{limite-monotone}{
  name={théorème de la limite monotone},
  description={Toute suite monotone bornée converge~; une suite croissante majorée tend vers son supremum, une suite décroissante minorée tend vers son infimum}
}

\newglossaryentry{sous-suite}{
  name={sous-suite},
  description={Suite extraite $(u_{\varphi(n)})$ où $\varphi:\mathbb{N}\to\mathbb{N}$ est strictement croissante}
}

\newglossaryentry{borne}{
  name={suite bornée},
  description={Suite $(u_n)$ telle que $\exists M>0,\;\forall n,\;|u_n|\leq M$}
}

\newglossaryentry{valeur-adherence}{
  name={valeur d'adhérence},
  description={Limite d'une sous-suite convergente extraite de $(u_n)$}
}

\newglossaryentry{compacite}{
  name={compacité séquentielle},
  description={Propriété d'un espace dans lequel toute suite admet une sous-suite convergente}
}

\newglossaryentry{continuite}{
  name={continuité},
  description={$f:I\to\mathbb{R}$ est continue en $a$ si $\lim_{x\to a}f(x)=f(a)$}
}

\newglossaryentry{tvi}{
  name={théorème des valeurs intermédiaires},
  description={Toute fonction continue sur un intervalle $[a,b]$ prend toutes les valeurs entre $f(a)$ et $f(b)$}
}

\newglossaryentry{connexite}{
  name={connexité par arcs},
  description={Propriété topologique des intervalles de $\mathbb{R}$~: deux points sont toujours reliés par un chemin continu}
}

\newglossaryentry{endomorphisme}{
  name={endomorphisme},
  description={Application linéaire d'un espace vectoriel dans lui-même}
}

\newglossaryentry{valeur-propre}{
  name={valeur propre},
  description={Scalaire $\lambda$ tel qu'il existe $v\neq 0$ avec $f(v)=\lambda v$}
}

\newglossaryentry{vecteur-propre}{
  name={vecteur propre},
  description={Vecteur non nul $v$ tel que $f(v)=\lambda v$ pour une valeur propre $\lambda$}
}

\newglossaryentry{diagonalisable}{
  name={diagonalisable},
  description={Endomorphisme admettant une base de vecteurs propres}
}

\newglossaryentry{autoadjoint}{
  name={autoadjoint (opérateur)},
  description={Endomorphisme $f$ d'un espace euclidien tel que $\langle f(x),y\rangle=\langle x,f(y)\rangle$ pour tous $x,y$}
}

\newglossaryentry{polynome-caracteristique}{
  name={polynôme caractéristique},
  description={Polynôme $\chi_f(\lambda)=\det(\lambda\,\mathrm{Id}-f)$ associé à un endomorphisme $f$}
}

\newglossaryentry{polynome-minimal}{
  name={polynôme minimal},
  description={Polynôme unitaire de plus petit degré qui annule un endomorphisme}
}

\newglossaryentry{wallis}{
  name={intégrales de Wallis},
  description={Intégrales $W_n=\int_0^{\pi/2}\sin^n(t)\,\mathrm{d}t$, intervenant notamment dans la formule de Stirling}
}

\newglossaryentry{stirling}{
  name={formule de Stirling},
  description={Équivalent asymptotique $n!\sim\sqrt{2\pi n}\,{(n/e)}^n$}
}

\newglossaryentry{espace-euclidien}{
  name={espace euclidien},
  description={Espace vectoriel réel de dimension finie muni d'un produit scalaire}
}

% --- Structures algébriques ---

\newglossaryentry{groupe}{
  name={groupe},
  description={Ensemble muni d'une loi associative, d'un neutre et d'un symétrique pour tout élément}
}

\newglossaryentry{groupe-abelien}{
  name={groupe abélien},
  description={Groupe dont la loi est commutative~; noté additivement $(G,+)$}
}

\newglossaryentry{sous-groupe-distingue}{
  name={sous-groupe distingué},
  description={Sous-groupe $H \trianglelefteq G$ stable par conjugaison~: $gHg^{-1}=H$ pour tout $g\in G$}
}

\newglossaryentry{groupe-quotient}{
  name={groupe quotient},
  description={Groupe $G/H$ des classes d'un sous-groupe distingué $H$, muni de $(aH)(bH)=abH$}
}

\newglossaryentry{ordre-element}{
  name={ordre d'un élément},
  description={Plus petit entier $k\geq 1$ tel que $g^k=e$ dans un groupe~; divise l'ordre du groupe (Lagrange)}
}

\newglossaryentry{groupe-cyclique}{
  name={groupe cyclique},
  description={Groupe engendré par un seul élément~; isomorphe à $\mathbb{Z}$ ou $\mathbb{Z}/n\mathbb{Z}$}
}

\newglossaryentry{ideal}{
  name={idéal},
  description={Sous-groupe additif $I$ d'un anneau $A$ stable par multiplication externe~: $AI\subseteq I$}
}

\newglossaryentry{ideal-premier}{
  name={idéal premier},
  description={Idéal propre $\mathfrak{p}$ tel que $ab\in\mathfrak{p}\Rightarrow a\in\mathfrak{p}$ ou $b\in\mathfrak{p}$~; équivalent à $A/\mathfrak{p}$ intègre}
}

\newglossaryentry{ideal-maximal}{
  name={idéal maximal},
  description={Idéal propre non contenu dans aucun idéal propre plus grand~; équivalent à $A/\mathfrak{m}$ corps}
}

\newglossaryentry{extension-corps}{
  name={extension de corps},
  description={Inclusion $K\subseteq L$ de corps~; le degré $[L:K]=\dim_K L$ mesure la taille de l'extension}
}

\newglossaryentry{polynome-minimal-ext}{
  name={polynôme minimal (extension)},
  description={Polynôme unitaire irréductible de plus petit degré annulant un élément algébrique $\alpha$ sur $K$}
}

\newglossaryentry{corps-finis}{
  name={corps fini},
  description={Corps à $q=p^n$ éléments ($p$ premier), unique à isomorphisme près, noté $\mathbb{F}_q$~; son groupe multiplicatif est cyclique}
}

% --- Exponentielle et logarithme matriciels ---

\newglossaryentry{euler-exp-limite}{
  name={exponentielle matricielle (définition limite)},
  description={Pour $X\in\mathcal{M}_n(\mathbb{K})$~: $e^X=\lim_{k\to\infty}{(I+X/k)}^k$~; la convergence est en $O(\|X\|^2 e^{\|X\|}/k)$ et généralise la limite des intérêts composés}
}

\newglossaryentry{trotter}{
  name={formule de Trotter},
  description={Pour $A,B\in\mathcal{M}_n(\mathbb{K})$ quelconques~: $e^{A+B}=\lim_{k\to\infty}{(e^{A/k}e^{B/k})}^k$~; repose sur Baker--Campbell--Hausdorff}
}

\newglossaryentry{jacobi-liouville}{
  name={formule de Jacobi--Liouville},
  description={Pour toute matrice $A$~: $\det(e^A)=e^{\mathrm{tr}(A)}$~; implique $\det(e^A)>0$ sur $\mathbb{R}$ et $\exp(\mathfrak{sl}_n)\subseteq\mathrm{SL}_n$}
}

\newglossaryentry{liouville-wronskien}{
  name={formule de Liouville (wronskien)},
  description={Pour $X'=A(t)X$, le wronskien vérifie $W(t)=W(0)\exp\!\int_0^t\mathrm{tr}(A(s))\,ds$~; généralise $\det(e^{tA})=e^{t\,\mathrm{tr}(A)}$}
}

\newglossaryentry{logarithme-matriciel}{
  name={logarithme matriciel},
  description={Antécédent $A$ de $M$ par l'exponentielle~: $e^A=M$~; pour un bloc de Jordan $J_m(\lambda)$ avec $\lambda\neq 0$, il vaut $\mu I_m + \sum_{k=1}^{m-1}\frac{{(-1)}^{k+1}}{k}{(N/\lambda)}^k$ où $e^\mu=\lambda$}
}

\newglossaryentry{surjectivite-exp}{
  name={surjectivité de l'exponentielle matricielle},
  description={$\exp:\mathcal{M}_n(\mathbb{C})\to\mathrm{GL}_n(\mathbb{C})$ est surjective (via le logarithme de Jordan)~; sur $\mathbb{R}$, l'image est contenue dans $\mathrm{GL}_n^+(\mathbb{R})$ mais l'inclusion est stricte}
}

% --- Polynômes classiques ---

\newglossaryentry{tchebychev}{
  name={polynôme de Tchebychev},
  description={Polynôme $T_n$ défini par $T_n(\cos\theta)=\cos(n\theta)$~; vérifie $T_{n+1}=2XT_n-T_{n-1}$, de coefficient dominant $2^{n-1}$ pour $n\geq 1$}
}

\newglossaryentry{tchebychev-approx}{
  name={approximation minimale (Tchebychev)},
  description={Parmi les polynômes unitaires de degré $n$, $\frac{1}{2^{n-1}}T_n$ minimise la norme uniforme sur $[-1,1]$, qui vaut $\frac{1}{2^{n-1}}$}
}

\newglossaryentry{noeuds-tchebychev}{
  name={nœuds de Tchebychev},
  description={Racines de $T_{n+1}$~: $x_k=\cos\!\bigl(\frac{(2k-1)\pi}{2(n+1)}\bigr)$, $k=1,\ldots,n+1$~; minimisent l'erreur d'interpolation sur $[-1,1]$}
}

% --- Combinatoire ---

\newglossaryentry{nombre-bell}{
  name={nombre de Bell},
  description={Entier $B_n$ comptant le nombre de partitions de l'ensemble $\{1,\ldots,n\}$~; $B_0=1$, $B_1=1$, $B_2=2$, $B_3=5$, $B_4=15$}
}

\newglossaryentry{partition-ensemble}{
  name={partition d'un ensemble},
  description={Famille de parties non vides, deux à deux disjointes, dont la réunion est l'ensemble tout entier}
}

\newglossaryentry{stirling2}{
  name={nombre de Stirling de deuxième espèce},
  description={Entier $S(n,k)$ comptant les partitions de $\{1,\ldots,n\}$ en exactement $k$ blocs non vides~; vérifie $S(n,k)=kS(n-1,k)+S(n-1,k-1)$ et $B_n=\sum_k S(n,k)$}
}

% --- Construction des nombres ---

\newglossaryentry{axiomes-peano}{
  name={axiomes de Peano},
  description={Système axiomatique caractérisant $\mathbb{N}$~: élément $0$, successeur injectif ne touchant pas $0$, et principe de récurrence}
}

\newglossaryentry{successeur}{
  name={successeur},
  description={Application $S:\mathbb{N}\to\mathbb{N}$ définie par $S(n)=n+1$, fondamentale dans les axiomes de Peano}
}

\newglossaryentry{bonne-fondation}{
  name={bonne fondation},
  description={Propriété de $\mathbb{N}$~: tout sous-ensemble non vide admet un plus petit élément}
}

\newglossaryentry{von-neumann}{
  name={construction de von Neumann},
  description={Modèle ensembliste de $\mathbb{N}$ dans ZFC~: $0=\emptyset$, $S(n)=n\cup\{n\}$}
}

\newglossaryentry{entiers-relatifs}{
  name={entiers relatifs},
  description={Ensemble $\mathbb{Z}=(\mathbb{N}\times\mathbb{N})/{\sim}$ où $(m,n)\sim(m',n')\Leftrightarrow m+n'=m'+n$, représentant la différence $m-n$}
}

\newglossaryentry{anneau-initial}{
  name={anneau initial},
  description={$\mathbb{Z}$ est l'anneau initial~: il existe un unique homomorphisme d'anneaux $\mathbb{Z}\to A$ pour tout anneau $A$}
}

\newglossaryentry{coupure-dedekind}{
  name={coupure de Dedekind},
  description={Partition $(\mathcal{L},\mathcal{U})$ de $\mathbb{Q}$ en un segment initial sans maximum et son complémentaire, modélisant un réel}
}

\newglossaryentry{suite-cauchy}{
  name={suite de Cauchy},
  description={Suite $(a_n)$ dans $\mathbb{Q}$ (ou $\mathbb{R}$) telle que $\forall\varepsilon>0,\;\exists N,\;\forall n,m\geq N,\;|a_n-a_m|<\varepsilon$}
}

\newglossaryentry{completude}{
  name={complétude},
  description={Propriété d'un espace métrique dans lequel toute suite de Cauchy converge~; $\mathbb{R}$ est la complétion de $\mathbb{Q}$}
}

\newglossaryentry{corps-alg-clos}{
  name={corps algébriquement clos},
  description={Corps dans lequel tout polynôme non constant admet une racine~; $\mathbb{C}$ est la clôture algébrique de $\mathbb{R}$}
}

\newglossaryentry{quaternions}{
  name={quaternions},
  description={Algèbre $\mathbb{H}=\{a+bi+cj+dk\mid a,b,c,d\in\mathbb{R}\}$ de dimension 4 sur $\mathbb{R}$, corps gauche non commutatif}
}

\newglossaryentry{corps-gauche}{
  name={corps gauche},
  description={Anneau non nul dans lequel tout élément non nul est inversible, sans supposer la commutativité du produit~; exemple~: $\mathbb{H}$}
}

\newglossaryentry{ZFC}{
  name={ZFC},
  description={Théorie des ensembles de Zermelo-Fraenkel avec axiome du choix, cadre standard de la mathématique moderne}
}
